%%%%%%%%%%  scan idx 66  |  folio 53  %%%%%%%%%%
\emph{(2.22) est exacte~; si $Z$ est homogène de codimension $i$ et si $f^{-1}(Z)$ n'a pas de composante excédentaire, on a~:}
\begin{equation}\tag{2.23}
 f^!(\gamma_Y(Z))\equiv\gamma_X(f^{-1}(Z)).\qquad\mathrm{mod.}\ \underline K(X)^{i+1}
\end{equation}
Démonstration analogue à celle de la proposition 2.7.

Signalons enfin qu'on verra au \S 4 que, \emph{lorsque $X$ et $Y$ sont quasi-projectives et non singulières, alors $f^!$ est compatible avec les filtrations de $\underline K(X)$ et de $\underline K(Y)$}.

Soient de nouveau $X$ et $Y$ deux espaces algébriques quelconques, et $f$ un morphisme \emph{propre} de $X$ dans $Y$. Pour tout faisceau algébrique cohérent $F$ sur $X$, le faisceau image directe $f_!(F)$ est algébrique cohérent, et il en est de même plus généralement des faisceaux $R^qf_!(F)$ (cf. [3]). Rappelons que, par définition, $R^qf_!(F)$ est le faisceau sur $Y$ associé au préfaisceau qui à l'ouvert $U$ fait correspondre le groupe $H^q(f^{-1}(U),F)$~; on a en particulier~:
\begin{equation}\tag{2.24}
 R^0f_!(F)=f_!(F),\qquad R^qf_!(F)=0\quad(q>\dim X),
\end{equation}
de plus~:
\begin{equation}\tag{2.25}
 \Gamma(U,R^qf_!(F))=H^q(f^{-1}(U),F)
\end{equation}
\emph{si l'ouvert $U$ est affine}. D'après le \S 2, on définit un homomorphisme de $\underline K(X)$ dans $\underline K(Y)$, noté encore $f_!$, par la formule
\begin{equation}\tag{2.26}
 f_!(\gamma_X(F))=\sum_{q\geq0}(-1)^q\gamma_Y(R^qf_!(F)).
\end{equation}
Ainsi $\underline K(X)$ devient un \emph{foncteur covariant en $X$}, relativement aux morphismes \emph{propres}. La relation $(fg)_!=f_!g_!$ (tout comme la relation $(fg)^!=g^!f^!$ qu'on a oubliée de signaler en son lieu) est un cas
parti%
